\documentclass[xcolor=table,aspectratio=169]{beamer}
\usepackage{beamerthemesplit}
\usepackage{wrapfig}
\usetheme{SPbGU}
\usepackage{pdfpages}
\usepackage{amsmath}
\usepackage{cmap}
\usepackage[T2A]{fontenc}
\usepackage[utf8]{inputenc}
\usepackage[english]{babel}
\usepackage{indentfirst}
\usepackage{amsmath}
\usepackage{tikz}
\usepackage{multirow}
\usepackage[noend]{algpseudocode}
\usepackage{algorithm}
\usepackage{algorithmicx}
\usepackage{fancyvrb}
\usetikzlibrary{shapes,arrows}
%usepackage{fancyvrb}
%\usepackage{minted}
%\usepackage{verbments}
\usepackage{fontawesome}


\beamertemplatenavigationsymbolsempty

\title[CFPQ: Obstacles on the Way to Adoption]{Context-Free Path Querying: Obstacles on the Way to Adoption}
\institute[JBR, SPbSU]{
JetBrains Research, Programming Languages and Tools Lab  \\
St. Petersburg State University
}

\author[Semyon Grigorev]{Semyon Grigorev}

\date{16.07.2021}

\begin{document}
{
\begin{frame}[fragile]
  \begin{tabular}{p{2.0cm} p{9.5cm} p{1cm}}
   \begin{center}
      \includegraphics[height=1.5cm]{pictures/jetbrainsResearch.pdf}
    \end{center}
    &
    \begin{center}
      LDBC TUC 
    \end{center}
    &
    \begin{center}
      \includegraphics[height=1.5cm]{pictures/SPbGU_Logo.png}
    \end{center}
  \end{tabular}
  \titlepage
\end{frame}
}

\begin{frame}[fragile]

  \frametitle{Who we are?}

  \begin{minipage}[t]{0.48\textwidth}
    \begin{itemize}
    \item !!!
    \end{itemize}    
  \end{minipage}
  \begin{minipage}[t]{0.48\textwidth}
    \begin{itemize}
    \item !!!
    \end{itemize}    
  \end{minipage}

\end{frame}

\begin{frame}[fragile]
 
  \frametitle{Formal language constarined path querying}

  \begin{minipage}[t]{0.58\textwidth}
    Navigational queries in edge-labelled graph
    \begin{itemize}
    \item $w(v_0 \xrightarrow{l_0} v_1 \xrightarrow{l_1} \ldots \xrightarrow{l_{k-1}} v_k) = l_0 l_1 \ldots l_{k-1}$
    \item $Q = \{(v_i,v_j) \mid \exists \pi = v_i \to \ldots \to v_j; w(\pi) \in \mathcal{L}\}$, where $\mathcal{L}$ --- formal language
    \end{itemize}    
  \end{minipage}
  \begin{minipage}[t]{0.38\textwidth}
    \begin{itemize}
    \item !!!
    \end{itemize}    
  \end{minipage}

\end{frame}


\begin{frame}[fragile]
  \frametitle{Applications of Context-Free Path Querying}
  \begin{minipage}[t]{0.48\textwidth}
    Static code analysis
  \end{minipage}
  \begin{minipage}[t]{0.48\textwidth}
    Static code analysis
  \end{minipage}
\end{frame}

\begin{frame}[fragile]
  \frametitle{Problems}
  \begin{itemize}    
    \item There is no unified infrastructure for solutions comparison
    \begin{itemize}
      \item Data is spread over articles in different communities
      \item There is a huge number of different subclasses of the problem
      \begin{itemize}
        \item all-pairs, single source, multiple source, \dots 
        \item reachability, single path, all path, \dots 
      \end{itemize}
      \item The first and only attempt to compare different algorithms: ``An Experimental Study of Context-Free Path Query Evaluation Methods''\footnote{Jochem Kuijpers, George Fletcher, Nikolay Yakovets, and Tobias Lindaaker. 2019}
      \pause
      \item Conclusion: ``We conclude that state of the art solutions are not able to cope with large graphs as found in practice.''
    \end{itemize}
    \pause
    \item There is no support in real-world graph database
    \begin{itemize}
      \item H. Miao and A. Deshpande: ``Though the problem has been first studied in our community [40], there is little follow up and support in the context of modern graph databases \dots''\footnote{H. Miao and A. Deshpande,  ``Understanding Data Science Lifecycle Provenance via Graph Segmentation and Summarization'', 2019}       
    \end{itemize}
    
  \end{itemize}
\end{frame}

\begin{frame}[fragile]
  \frametitle{Our Results}

  \begin{itemize}
    \item[\faCheck] Collection of linear algebra based algorithms for CFPQ
    \begin{itemize}
      \item SuiteSparse is utilized for sparse linear algebra subroutines
      \item Published: \url{https://github.com/JetBrains-Research/CFPQ_PyAlgo}
    \end{itemize}
    \pause
    % Definitely, not all abilities of Cypher language is supported, but ...
    \item[\faCheck] Full-stack support of CFPQ\footnote{``Multiple-Source Context-Free Path Querying in Terms of Linear Algebra'', Arseniy Terekhov et al, 2019}    
    \begin{itemize}
      \item On the top of RedisGraph
      \item openCypher extended to support CFPQ
    \end{itemize} 
    \pause 
    \item[\faGears] Collecting of the dataset for CFPQ benchmarking is started
    \begin{itemize}
      \item Synthetic graphs
      \item Real-world graphs
      \begin{itemize}
        \item Static code analysis
        \item Biological data analysis
      \end{itemize}
      \item Published: \url{https://github.com/JetBrains-Research/CFPQ_Data} 
    \end{itemize}     
  \end{itemize}

\end{frame}

\begin{frame}[fragile]
  \frametitle{Our Results Evaluation}
  \begin{minipage}[t]{0.48\textwidth}
    \begin{itemize}
      \item Reachability queries
      \begin{itemize}
        \item \textit{geospecies} --- biological data, folklore
        \item \textit{mem} --- Andersen points-to analysis
      \end{itemize}
      \item Time in seconds
    \end{itemize}
  \end{minipage}
  \begin{minipage}[t]{0.48\textwidth}
    \begin{itemize}
      \item GPGPU: Geforce GTX 1070, 1.5GHz, 8Gb RAM, 1920 CUDA cores 
      \item CPU: Intel core i7-6700 CPU, 3.4GHz, DDR4 64Gb RAM
    \end{itemize}
  \end{minipage}
\begin{center}
  \rowcolors{2}{black!2}{black!10}
  \begin{tabular}{| c | c | c || c | c | c | c |}
    \hline
    Graph      & \#V      & \#E       & Neo4j\footnote{!!!} & RedisGraph\footnote{!!!} & Lin.al. CPU\footnote{!!!} & Lin.al. GPGPU\footnote{!!!} \\
    \hline
    \hline
    geospecies & 450 609  & 2 311 461 & 6000.0              & 80.1        & 7.1           & 0.8               \\ 
    Mem 1      &          &           & n.a.                & \faGears    &               &                   \\
    mem 2      &          &           & n.a.                & \faGears    &               &                   \\
    mem 3      &          &           & n.a.                & \faGears    &               &                   \\
    \hline
  \end{tabular}  
\end{center}
\end{frame}


\begin{frame}[fragile]

  \frametitle{Ongoing research}
\begin{itemize}
      \item[\faGears] Benchmarking of linear algebra based algorithms
      \begin{itemize}
        \item Comparison of different algorithms for different query semantics
        \item Investigation of scalability on multicore machines
        \item Estimation of performance on GPGPU
      \end{itemize}
      \item[\faGears] GLL-based CFPQ algorithm for Neo4j
      \pause
      \item[\faHourglassHalf] Semantics of openCypher in terms of linear algebra (in Coq)
      \item[\faHourglassHalf] Multiple context-free languages for constraints      
\end{itemize}
\end{frame}

\begin{frame}[fragile]
  \frametitle{Topics to dicuss}
  \begin{itemize}
    \item Benchmarks. 
    \begin{itemize}
      \item Algorithms
      \item Query language. Semantics 
      \item Local and Global queries. Can start from RPQ.
      \item Queries: templates and real-world queries.
    \end{itemize}
    \item Competition on language constrained path querying.   
    \begin{itemize}
      \item Part of existing competition for graph processing systems.
      \item Involve static analysis community. 
    \end{itemize}
    \item Graph database support
    \begin{itemize}
      \item Algorithms
      \item Query language. Semantics 
    \end{itemize}
  \end{itemize}
\end{frame}

\end{document}

